\section{Finite \texorpdfstring{$\Ein$}{Ein} witnesses and
one-variable Euler traces}
\label{sec:finite-witness-zero-traces}

The pure-exponential base-change theorem also controls finite
representations of constants by algebraic-point \(\Ein\)-values.  We apply
that observation first to \(\Ein(z)-\gamma\), and then to the rank-one
witness family
\(n\Ein(z)-\Ein(z^n)-(n-1)\gamma\).  No appeal to Picard's theorem or to a
numerical zero count is needed.

\subsection{Rigidity of finite algebraic-point witnesses}

Put
\[
 \mathscr V=\operatorname{span}_{\Qbar}
 \{\Ein(\alpha):\alpha\in\Qbar^\times\}\subset\C,
\]
and let \(\Eexp^{\mathrm{alg}}\) denote the relative algebraic closure of
\(\Eexp\) in \(\C\).

\begin{proposition}[Finite-witness rigidity]
\label{prop:finite-ein-witness-rigidity}
The linear map
\[
 \bigoplus_{\alpha\in\Qbar^\times}\Qbar e_\alpha
 \longrightarrow\C,
 \qquad
 \sum_\alpha c_\alpha e_\alpha\longmapsto
 \sum_\alpha c_\alpha\Ein(\alpha),
\]
where the direct sum means finite support, is injective.  Moreover
\begin{equation}\label{eq:V-Eexp-intersection}
 \mathscr V\cap\Eexp^{\mathrm{alg}}=\{0\}.
\end{equation}
Consequently, if \(\gamma\in\mathscr V\), then its finite
\(\Ein\)-representation is unique and \(\gamma\) is transcendental over
\(\Eexp\).
\end{proposition}

\begin{proof}
Every finite family of distinct nonzero algebraic-point values
\(\Ein(\alpha)\) is algebraically independent over \(\Eexp\) by
Theorem~\ref{thm:pure-exp-base}; it is therefore linearly independent.
This proves injectivity.  A nonzero linear form in algebraically independent
indeterminates is itself transcendental over the base field: extend that
linear form to an invertible linear coordinate system.  Hence every nonzero
element of \(\mathscr V\) is transcendental over \(\Eexp\), proving
\eqref{eq:V-Eexp-intersection} and the final assertion.
\end{proof}

Let \(q_*\in(0,1)\) be the unique real solution of
\begin{equation}\label{eq:qstar-ein-level}
 \Ein(q_*)=\gamma.
\end{equation}
Indeed, \(\Ein'(x)=(1-e^{-x})/x>0\) on \(\mathbb R\), after filling the
removable value at \(0\), while
\(\Ein(0)-\gamma=-\gamma<0\) and
\(\Ein(1)-\gamma=\delta/e>0\).

\begin{theorem}[The one-variable Euler-level dichotomy]
\label{thm:qstar-algebraic-dichotomy}
Exactly one of the following alternatives holds.
\begin{enumerate}[label=\textup{(\alph*)}]
 \item Every zero of \(\Ein(z)-\gamma\) is transcendental.
 \item The number \(q_*\) is algebraic and is the unique algebraic zero of
       \(\Ein(z)-\gamma\).  In this case \(\gamma\) and \(\delta\) are
       algebraically independent over \(\Eexp\).
\end{enumerate}
Under alternative \textup{(b)}, the six-constant field
\[
 \Qbar\bigl(\gamma,\delta,E_1(1),\Ei(1),\Ci(1),\Chi(1)\bigr)
\]
has transcendence degree \(5\), and its complete relation ideal is
generated by
\begin{equation}\label{eq:qstar-six-relation}
 \Ei(1)-E_1(1)-2\Chi(1).
\end{equation}
\end{theorem}

\begin{proof}
Two distinct algebraic zeros \(a,b\) would give
\(\Ein(a)=\Ein(b)\), contradicting
Theorem~\ref{thm:pure-exp-base}.  A nonreal algebraic zero is also
impossible, since its distinct conjugate would be a second algebraic zero.
Thus the only possible algebraic zero is the unique real zero \(q_*\).

Assume that \(q_*\) is algebraic and put \(Y_a=\Ein(a)\).  The values
\(Y_{q_*}\) and \(Y_1\) are independent over \(\Eexp\), and
\[
 \gamma=Y_{q_*},\qquad
 \frac{\delta}{e}=Y_1-Y_{q_*}.
\]
This invertible linear change proves the asserted independence of
\(\gamma,\delta\) over \(\Eexp\).

For the last assertion, put
\[
 A=\Ein(1),\qquad B=\Ein(-1),\qquad
 C=\frac{\Ein(i)+\Ein(-i)}2.
\]
The five arguments \(q_*,1,-1,i,-i\) are distinct.  Pure-exponential
base change shows that \(\gamma=\Ein(q_*),A,B,C\) are independent over
\(\Eexp\); together with the transcendence of \(e\), the five numbers
\(e,\gamma,A,B,C\) are algebraically independent over \(\Qbar\).
The six-constant field is \(\Qbar(e,\gamma,A,B,C)\).  Identity
\eqref{eq:qstar-six-relation} is therefore a height-one prime relation and
generates the complete kernel.
\end{proof}

\subsection{The rank-one witness family}

For \(n\ge2\), define
\begin{equation}\label{eq:Gn-witness-family}
 G_n(z)=n\Ein(z)-\Ein(z^n)-(n-1)\gamma.
\end{equation}

\begin{proposition}[Real-zero classification for \(G_n\)]
\label{prop:Gn-real-zero-classification}
For every \(n\ge2\), the function \(G_n\) has exactly one positive real zero
\(\beta_n\in(0,1)\).  If \(n\) is even, it has no negative real zero.  If
\(n\) is odd, it has exactly one negative real zero
\(\lambda_n\in(-\infty,-1)\).  All these real zeros are simple.
\end{proposition}

\begin{proof}
For \(x\ne0\), direct differentiation gives
\begin{equation}\label{eq:Gn-real-derivative}
 G_n'(x)=\frac n{x}\bigl(e^{-x^n}-e^{-x}\bigr).
\end{equation}
Thus \(G_n\) is strictly increasing on \((0,1)\).  Since
\[
 G_n(0)=-(n-1)\gamma<0,
 \qquad
 G_n(1)=\frac{(n-1)\delta}{e}>0,
\]
there is exactly one zero \(\beta_n\) in that interval.  For \(x>1\),
\[
 G_n(x)=nE_1(x)-E_1(x^n)>(n-1)E_1(x)>0,
\]
so there are no further positive zeros.

If \(n\) is even and \(x<0\), then \(x^n>0>x\); hence the numerator in
\eqref{eq:Gn-real-derivative} is negative and \(G_n'(x)>0\).  Moreover
\(G_n(x)\to-\infty\) as \(x\to-\infty\), while
\(G_n(0)=-(n-1)\gamma<0\).  Strict monotonicity therefore excludes a
negative zero.  The asserted limit also follows directly, on writing
\(x=-t\), from
\[
 G_n(-t)=-n\Ei(t)-E_1(t^n)\qquad(t>0).
\]

If \(n\) is odd, \eqref{eq:Gn-real-derivative} shows that \(G_n\) is
strictly decreasing on \((-\infty,-1)\) and strictly increasing on
\((-1,0)\).  One has
 \begin{align*}
 G_n(x)&\longrightarrow+\infty\quad(x\to-\infty),\\
 G_n(-1)&=-(n-1)\Ei(1)<0,\qquad
 G_n(0)=-(n-1)\gamma<0.
 \end{align*}
Here the first limit follows from
\[
 G_n(-t)=\Ei(t^n)-n\Ei(t)\qquad(t>0)
\]
and the standard positive-axis asymptotic for \(\Ei\).
This gives exactly one zero in \((-\infty,-1)\) and none in \((-1,0)\).
Formula \eqref{eq:Gn-real-derivative} is nonzero at each listed root, so
they are simple.
\end{proof}

\begin{theorem}[Algebraic-zero rigidity for \(G_n\)]
\label{thm:Gn-algebraic-zero-rigidity}
For every \(n\ge2\), the function \(G_n\) has at most one algebraic zero,
and every algebraic zero is real.  Consequently:
\begin{enumerate}[label=\textup{(\roman*)}]
 \item if \(n\) is even, either every zero is transcendental or
       \(\beta_n\) is the unique algebraic zero;
 \item if \(n\) is odd, every nonreal zero is transcendental and at most
       one of \(\beta_n,\lambda_n\) is algebraic.
\end{enumerate}
If \(\alpha\) is an algebraic zero and
\(D_\alpha=\Ein(\alpha)-\Ein(\alpha^n)\), then
\begin{equation}\label{eq:Gn-witness-amplification}
 \gamma,\quad\delta,\quad D_\alpha
 \quad\text{are algebraically independent over }\Eexp.
\end{equation}
The six-constant field then has transcendence degree \(5\), with complete
relation ideal generated by \eqref{eq:qstar-six-relation}.
\end{theorem}

\begin{proof}
Suppose that distinct algebraic numbers \(\alpha,\beta\) were zeros.  They
are nonzero, and subtraction gives
\[
 n\Ein(\alpha)-\Ein(\alpha^n)
 =n\Ein(\beta)-\Ein(\beta^n).
\]
After collecting equal arguments, this is a linear relation among distinct
algebraic-point \(\Ein\)-values.  Its coefficient at
\(\Ein(\alpha)\) is
\[
 n-\mathbf1_{\alpha^n=\alpha}
   +\mathbf1_{\beta^n=\alpha}\ge n-1>0,
\]
because \(\beta\ne\alpha\).  This contradicts
Proposition~\ref{prop:finite-ein-witness-rigidity}.  Hence there is at most
one algebraic zero.  Since \(G_n\) has real coefficients, a nonreal
algebraic zero and its conjugate would give two, so every algebraic zero is
real.  Proposition~\ref{prop:Gn-real-zero-classification} now gives the two
alternatives.

For either possible algebraic real zero, the three arguments
\(\alpha,\alpha^n,1\) are distinct.  The corresponding three
\(\Ein\)-values are independent over \(\Eexp\), while
\[
 \gamma=\frac{n\Ein(\alpha)-\Ein(\alpha^n)}{n-1},\qquad
 \frac\delta e=\Ein(1)-\gamma,\qquad
 D_\alpha=\Ein(\alpha)-\Ein(\alpha^n).
\]
This is an invertible linear change of three coordinates, proving
\eqref{eq:Gn-witness-amplification}.  Moreover
\(\{\alpha,\alpha^n\}\) is disjoint from \(\{1,-1,i,-i\}\), so the final
six-constant assertion follows exactly as in
Theorem~\ref{thm:qstar-algebraic-dichotomy}.
\end{proof}

\subsection{Two exact traces of the one-variable Euler divisor}

Let \(Z_\gamma\) be the multiset of zeros of
\(F_\gamma(z)=\Ein(z)-\gamma\), and put
\begin{equation}\label{eq:one-variable-traces}
 S_m=\sum_{\rho\in Z_\gamma}\rho^{-m}\qquad(m\ge2).
\end{equation}
The sums converge absolutely because \(F_\gamma\) is entire of order one.

\begin{theorem}[Trace recovery for \(\Ein-\gamma\)]
\label{thm:ein-gamma-trace-recovery}
One has
\begin{align}
 S_2&=\frac1{\gamma^2}-\frac1{2\gamma},
 \label{eq:ein-gamma-S2}\\
 S_3&=\frac1{\gamma^3}-\frac3{4\gamma^2}+\frac1{6\gamma}.
 \label{eq:ein-gamma-S3}
\end{align}
Conversely,
\begin{equation}\label{eq:ein-gamma-trace-inversion}
 \boxed{\displaystyle
 \gamma=\frac{24S_2+1}{24S_3+6S_2}.}
\end{equation}
In particular
\[
 \Qbar(S_2,S_3)=\Qbar(\gamma),
\]
and the two traces satisfy
\begin{equation}\label{eq:S2-S3-curve}
 144S_2^3+21S_2^2+S_2-144S_3^2+3S_3=0.
\end{equation}
Moreover \(S_2\in\Qbar\) if and only if \(\gamma\in\Qbar\), and the same
equivalence holds with \(S_3\) in place of \(S_2\).
\end{theorem}

\begin{proof}
A genus-one Hadamard factorization of \(F_\gamma\) gives, for \(m\ge2\),
the inverse power sums from the Taylor coefficients of
\(\log(F_\gamma(z)/F_\gamma(0))\).  The origin jet is
\[
 F_\gamma(0)=-\gamma,\quad
 F_\gamma'(0)=1,\quad
 F_\gamma''(0)=-\frac12,\quad
 F_\gamma'''(0)=\frac13.
\]
Hence
\[
 S_2=-\left(\log F_\gamma\right)''(0),
 \qquad
 S_3=-\frac12\left(\log F_\gamma\right)'''(0),
\]
which gives \eqref{eq:ein-gamma-S2} and
\eqref{eq:ein-gamma-S3}.  With \(x=1/\gamma\), these read
\[
 S_2=x^2-\frac x2,
 \qquad
 S_3=x^3-\frac34x^2+\frac x6.
\]
Eliminating \(x^2\) gives
\[
 x\left(S_2+\frac1{24}\right)=S_3+\frac{S_2}{4},
\]
which proves \eqref{eq:ein-gamma-trace-inversion}; its denominator is
nonzero because \(x=1/\gamma>1\).  A resultant calculation, or direct
substitution, gives \eqref{eq:S2-S3-curve}.  Finally, each \(S_j\) is a
nonconstant polynomial in \(x\) over \(\Qbar\), so algebraicity of either
trace forces \(x\), hence \(\gamma\), to be algebraic.  The converse is
immediate.
\end{proof}

\begin{remark}[All higher traces]
The same calculation gives
\[
 S_m=-m[z^m]\log\left(
  1+\frac1\gamma\sum_{r\ge1}
       \frac{(-1)^r z^r}{r\,r!}\right)
 \qquad(m\ge2).
\]
Thus every higher inverse trace is a rational polynomial in \(1/\gamma\).
Unlike the phase recovery from the zero set of the \(\gamma\)-free
reciprocal function, this is a spectral reparametrization of a
\(\gamma\)-dependent level.
\end{remark}
